\documentclass[10pt,a4paper]{amsart}
\usepackage[margin=19mm]{geometry}
\usepackage{amsmath,amssymb,mathtools}
\usepackage[scr=rsfs,cal=euler]{mathalfa}
\usepackage{xfrac}
\usepackage[unicode,hidelinks,bookmarks=false]{hyperref}
\newcommand{\ie}{i.e.\ }
\newcommand{\cf}{cf.\ }
\newcommand{\resp}{respectively\ }
\newcommand{\Subsection}{\subsection}
\providecommand{\nobreakdash}{\nobreak\hbox{-}}
\setlength{\emergencystretch}{2em}
\multlinegap0pt
\usepackage{amssymb}
\usepackage{enumerate}
\newtheorem{lemma}{Lemma}
\DeclareMathOperator*{\essinf}{ess\,inf}
\DeclareMathOperator*{\esssup}{ess\,sup}
\DeclareMathOperator{\osc}{osc}
\title{Bourgain, Brezis and Mironescu theorem for fractional Sobolev spaces with variable exponents}
\author{Minhyun Kim}
\date{Source: arXiv:2210.00574v1 (CC BY 4.0)}
\begin{document}
\maketitle

\noindent\emph{Source and licence.} Bourgain, Brezis and Mironescu theorem for fractional Sobolev spaces with variable exponents. Version arXiv:2210.00574v1 (CC BY 4.0), distributed by arXiv under the Creative Commons Attribution 4.0 International licence, http://creativecommons.org/licenses/by/4.0/, as recorded in the arXiv licence metadata for this exact version and shown on the abstract page https://arxiv.org/abs/2210.00574v1. Full licence notice: \texttt{licenses/arxiv-2210.00574-CC-BY-4.0.txt}.\par\smallskip

\noindent\textit{Editorial scope.} Counted unit: the article's lemma lem:liminf (source0.tex lines 211--217). The article's complete original proof of it is delivered with it (source0.tex lines 219--312, one \texttt{proof} environment of 94 lines). No mathematics is written, repaired or re-derived: the statement and the proof are the article's own lines, in the article's own notation, and no retained line is substituted or re-flowed.  Retained uncounted as the source-local context the proof needs: lines 83--85; lines 90--92.  Nothing was omitted from any retained span, so the package carries no omission record.  No retained line contains a citation, so the wrapper carries no bibliography and no bibliography entry.  No retained line contains a figure environment, an image-inclusion command or drawing code, so no image is delivered and the package declares no asset dependency.  Editorial formatting only, in addition: an AMS \texttt{amsart} preamble that restates the article's own declarations and macros used by the retained lines, namely the environments \texttt{lemma} on separate counters, together with the standard packages those lines need.  The retained environments print in this document as Lemma 1 in order of appearance, whereas the article prints the same items with the numbering of its own full document.  The complete native source of the article as distributed in the arXiv e-print for this exact version is bundled unchanged as \texttt{sources/131514/source0.tex}.
\begin{equation} \label{eq:p-bound}
1 < p_{-} := \essinf_{x, y \in \Omega} p(x, y) \leq \esssup_{x, y \in \Omega} p(x, y) =: p_{+} < +\infty,
\end{equation}

\begin{equation} \label{eq:p-log-Holder}
r^{-\osc_{B_{r}(x)} p(x, \cdot)} \leq L \quad \text{for all } x \in \Omega \text{ and } 0 < r < \min \lbrace \mathrm{dist}(x, \partial \Omega), 1 \rbrace
\end{equation}

\begin{lemma} \label{lem:liminf}
Let $\Omega \subset \mathbb{R}^{n}$ be open and assume that $p$ satisfies \eqref{eq:p-bound} and \eqref{eq:p-log-Holder}. Then
\begin{equation} \label{eq:liminf}
\int_{\Omega} K_{n, \bar{p}(x)} |\nabla u(x)|^{\bar{p}(x)} \,\mathrm{d}x \leq \liminf_{s \nearrow 1} s(1-s) \int_{\Omega} \int_{\Omega} \frac{|u(x)-u(y)|^{p(x, y)}}{|x-y|^{n+sp(x, y)}} \,\mathrm{d}y \,\mathrm{d}x
\end{equation}
for any $u \in C^{2}(\overline{\Omega})$.
\end{lemma}

\begin{proof}
Let us fix $\varepsilon > 0$, $\theta_{0} \in (0,1)$, $R > 0$ and $x \in \Omega \cap B_{R}$. Since $u \in C^{2}(\overline{\Omega})$, there exists a constant $C > 0$, depending on $R$ and $\|u\|_{C^{2}(\overline{\Omega})}$, such that
\begin{equation} \label{eq:second-difference}
|u(x+h)-u(x)-\nabla u(x) \cdot h| \leq C |h|^{2}
\end{equation}
for all $|h| < \mathrm{dist}(x, \partial \Omega)$. Since $p(x, \cdot)$ is continuous at $x$, we find $r_{1} \in (0, \mathrm{dist}(x, \partial \Omega))$, depending on $\varepsilon$, $x$, $|\nabla u(x)|$ and $p$, such that
\begin{equation} \label{eq:r1}
(1-\varepsilon) |\nabla u(x)|^{p(x, x)} \leq |\nabla u(x)|^{p(x, y)} \leq (1+\varepsilon) |\nabla u(x)|^{p(x, x)}
\end{equation}
holds for all $y \in B_{r_{1}}(x)$. Moreover, by continuity of $p(x, \cdot)$ at $x$ again, there exists $r_{2} \in (0, \mathrm{dist}(x, \partial \Omega))$, depending on $\varepsilon$, $x$ and $\theta_0$, such that
\begin{equation} \label{eq:r2}
(1-\varepsilon) \theta^{p(x, x)} \leq \theta^{p(x, y)} \leq (1+ \varepsilon) \theta^{p(x, x)}
\end{equation}
holds for all $y \in B_{r_{2}}(x)$ and $\theta \in [\theta_{0}, 1]$. We set $r_{0} = \min\lbrace r_{1}, r_{2}, 1 \rbrace$ and let $r \in (0, r_{0})$.

For $x \in \Omega \cap B_{R}$ and $h \in B_{r}$, we have from \eqref{eq:second-difference}
\begin{equation*}
|\nabla u(x) \cdot h| \leq |u(x+h)-u(x)| + C |h|^2.
\end{equation*}
Using an inequality
\begin{equation} \label{eq:alg-ineq}
(a+b)^{q} \leq (1+\varepsilon)a^{q} + C(q, \varepsilon) b^{q}, \quad a, b \geq 0, q > 1,
\end{equation}
we obtain
\begin{equation*}
|\nabla u(x) \cdot h|^{p(x, x+h)} \leq (1+\varepsilon) |u(x+h)-u(x)|^{p(x, x+h)} + C |h|^{2p(x, x+h)}
\end{equation*}
for some $C = C(R, \|u\|_{C^{2}(\overline{\Omega})}, p_{+}, p_{-}, \varepsilon) > 0$. Multiplying by $|h|^{-n-sp(x, x+h)}$ and integrating over $B_{r}$ with respect to $h$, we have
\begin{equation} \label{eq:liminf-I}
\begin{split}
I
&:= \int_{B_{r}} \frac{|\nabla u(x) \cdot h|^{p(x, x+h)}}{|h|^{n+sp(x,x+h)}} \,\mathrm{d}h \\
&\leq (1+\varepsilon) \int_{B_{r}} \frac{|u(x+h)-u(x)|^{p(x, x+h)}}{|h|^{n+sp(x, x+h)}} \,\mathrm{d}h + C \int_{B_{r}} \frac{|h|^{2p(x, x+h)}}{|h|^{n+sp(x, x+h)}} \,\mathrm{d}h
\end{split}
\end{equation}
for each $x \in \Omega \cap B_{R}$.

Let us estimate $I$ from below. By using the spherical coordinates, we write
\begin{equation*}
I = \int_{0}^{r} \int_{\mathbb{S}^{n-1}} |\nabla u(x) \cdot \omega|^{p(x, x+\rho \omega)} \rho^{-1+(1-s)p(x, x+\rho \omega)} \,\mathrm{d}\mathcal{H}^{n-1}(\omega) \,\mathrm{d}\rho.
\end{equation*}
We use the change of variables $\omega = A\xi$, where $A \in O(n)$ is a rotation satisfying $\nabla u(x) \cdot \omega = |\nabla u(x)| \xi_{n}$, so that we have
\begin{equation} \label{eq:I}
I = \int_{0}^{r} \int_{\mathbb{S}^{n-1}} |\nabla u(x)|^{p(x, x+\rho A\xi)} |\xi_{n}|^{p(x, x+\rho A \xi)} \rho^{-1+(1-s)p(x, x+\rho A\xi)} \,\mathrm{d}\mathcal{H}^{n-1}(\xi) \,\mathrm{d}\rho.
\end{equation}
Using \eqref{eq:r1} and \eqref{eq:r2}, we obtain
\begin{equation*}
I \geq (1-\varepsilon)^2 \left( \int_{0}^{r} \int_{\lbrace |\xi_{n}| \geq \theta_{0} \rbrace} |\xi_{n}|^{\bar{p}(x)} \rho^{-1+(1-s)p(x, x+\rho A\xi)} \,\mathrm{d}\mathcal{H}^{n-1}(\xi) \,\mathrm{d}\rho \right) |\nabla u(x)|^{\bar{p}(x)}.
\end{equation*}
Moreover, since the assumption \eqref{eq:p-log-Holder} yields
\begin{equation} \label{eq:log-Holder}
L^{-1} \leq \min\lbrace 1, \rho^{\osc_{B_{\rho}(x)}p(x, \cdot)} \rbrace \leq \rho^{p(x, x+\rho A\xi)-\bar{p}(x)} \leq \max\lbrace 1, \rho^{-\osc_{B_{\rho}(x)}p(x, \cdot)}\rbrace \leq L,
\end{equation}
$I$ can be further estimated as
\begin{equation*}
\begin{split}
I
&\geq \frac{(1-\varepsilon)^2}{L^{1-s}} \left( \int_{0}^{r} \int_{\lbrace |\xi_{n}| \geq \theta_{0} \rbrace} |\xi_{n}|^{\bar{p}(x)} \rho^{-1+(1-s) \bar{p}(x)} \,\mathrm{d}\mathcal{H}^{n-1}(\xi) \,\mathrm{d}\rho \right) |\nabla u(x)|^{\bar{p}(x)} \\
&\geq \frac{(1-\varepsilon)^2}{(1-s)L^{1-s}} \left( \int_{\lbrace |\xi_{n}| \geq \theta_{0} \rbrace} |\xi_{n}|^{\bar{p}(x)} \,\mathrm{d}\mathcal{H}^{n-1}(\xi) \right) \frac{r^{(1-s)\bar{p}(x)}}{\bar{p}(x)} |\nabla u(x)|^{\bar{p}(x)}.
\end{split}
\end{equation*}
Since
\begin{equation*}
\begin{split}
\int_{\lbrace |\xi_{n}| \geq \theta_{0} \rbrace} |\xi_{n}|^{\bar{p}(x)} \,\mathrm{d}\mathcal{H}^{n-1}(\xi)
&\geq \int_{\mathbb{S}^{n-1}} |\xi_{n}|^{\bar{p}(x)} \,\mathrm{d}\mathcal{H}^{n-1}(\xi) - \int_{\lbrace |\xi_{n}| < \theta_{0} \rbrace} \theta_{0}^{p_{-}} \,\mathrm{d}\mathcal{H}^{n-1}(\xi) \\
&\geq \bar{p}(x) K_{n, \bar{p}(x)} - \theta_{0} |\mathbb{S}^{n-1}|,
\end{split}
\end{equation*}
we arrive at
\begin{equation} \label{eq:liminf-I-lower}
I \geq \frac{(1-\varepsilon)^2}{(1-s)L^{1-s}} \left( \bar{p}(x) K_{n, \bar{p}(x)} - \theta_0 |\mathbb{S}^{n-1}| \right) \frac{r^{(1-s)\bar{p}(x)}}{\bar{p}(x)} |\nabla u(x)|^{\bar{p}(x)}.
\end{equation}

For the last term on the right-hand side of \eqref{eq:liminf-I}, we use the assumption \eqref{eq:p-bound} to have
\begin{equation} \label{eq:liminf-I-upper}
\int_{B_{r}} |h|^{-n+(2-s)p(x, x+h)} \,\mathrm{d}h \leq \int_{B_{r}} |h|^{-n+(2-s)p_{-}} \,\mathrm{d}h \leq |\mathbb{S}^{n-1}| \frac{r^{(2-s)p_{-}}}{(2-s)p_{-}}.
\end{equation}
Therefore, combining \eqref{eq:liminf-I}, \eqref{eq:liminf-I-lower}, \eqref{eq:liminf-I-upper} and integrating over $\Omega \cap B_{R}$ yield
\begin{equation*}
\begin{split}
&(1-\varepsilon)^{2} \frac{s}{L^{1-s}} \int_{\Omega \cap B_{R}} \left( \bar{p}(x) K_{n, \bar{p}(x)} - \theta_{0} |\mathbb{S}^{n-1}| \right) \frac{r^{(1-s)\bar{p}(x)}}{\bar{p}(x)} |\nabla u(x)|^{\bar{p}(x)} \,\mathrm{d}x \\
&\leq (1+\varepsilon) s(1-s) \int_{\Omega} \int_{\Omega} \frac{|u(x)-u(y)|^{p(x, y)}}{|x-y|^{n+sp(x, y)}} \,\mathrm{d}y \,\mathrm{d}x + C \frac{s(1-s)}{2-s} r^{(2-s)p_{-}}
\end{split}
\end{equation*}
for some $C > 0$ independent of $s$. We take $\liminf_{s \nearrow 1}$ on both sides to deduce
\begin{equation*}
\begin{split}
&(1-\varepsilon)^{2} \int_{\Omega \cap B_{R}} \left( \bar{p}(x) K_{n, \bar{p}(x)} - \theta_{0} |\mathbb{S}^{n-1}| \right) \frac{1}{\bar{p}(x)} |\nabla u(x)|^{\bar{p}(x)} \,\mathrm{d}x \\
&\leq (1+\varepsilon) \liminf_{s \nearrow 1} s(1-s) \int_{\Omega} \int_{\Omega} \frac{|u(x)-u(y)|^{p(x, y)}}{|x-y|^{n+sp(x, y)}} \,\mathrm{d}y \,\mathrm{d}x.
\end{split}
\end{equation*}
Since $\varepsilon$, $\theta_0$ and $R$ are arbitrarily chosen, we conclude \eqref{eq:liminf}.
\end{proof}

\end{document}
